% !TeX program = xelatex
% Model specification, revised 2026-10-24. Compile three times with XeLaTeX.
\documentclass[11pt,a4paper]{article}
\usepackage[margin=2.35cm,headheight=15pt]{geometry}
\usepackage{amsmath,amssymb,bm,mathtools}
\usepackage{booktabs,tabularx,longtable,array}
\usepackage[no-config]{fontspec}
\usepackage{needspace,enumitem,fancyhdr}
\usepackage[unicode,hidelinks]{hyperref}
\defaultfontfeatures{Ligatures=Common}
\defaultfontfeatures[\rmfamily,\sffamily]{Ligatures=Common,Mapping=}
\setmainfont[AutoFakeSlant=0.2]{Noto Serif CJK SC}
\setsansfont{Noto Sans CJK SC}
\setmonofont{Noto Sans Mono CJK SC}
\XeTeXlinebreaklocale "zh"
\XeTeXlinebreakskip=0pt plus 1pt minus 0.1pt
\setlength{\parindent}{2em}
\setlength{\parskip}{0.35em}
\setlength{\emergencystretch}{3em}
\setlength{\LTpre}{0.6em}
\setlength{\LTpost}{0.6em}
\renewcommand{\arraystretch}{1.22}
\setlist{itemsep=0.2em,topsep=0.4em,leftmargin=2em}
\allowdisplaybreaks[2]
\setcounter{tocdepth}{2}
\setcounter{secnumdepth}{3}
\numberwithin{equation}{section}
\renewcommand{\contentsname}{目录}
\renewcommand{\tablename}{表}
\pagestyle{fancy}
\fancyhf{}
\fancyhead[L]{\small Model 1023：学习与决策闭环}
\fancyhead[R]{\small 2026-10-24}
\fancyfoot[C]{\thepage}
\newcolumntype{L}[1]{>{\raggedright\arraybackslash}p{#1}}
\newcolumntype{Y}{>{\raggedright\arraybackslash}X}
\newcommand{\ind}{\mathbb{I}}
\newcommand{\norm}[1]{\left\lVert#1\right\rVert}
\newcommand{\clip}{\operatorname{clip}}
\newcommand{\logit}{\operatorname{logit}}
\newcommand{\sigmoid}{\operatorname{sigmoid}}
\newcommand{\A}{\mathcal{A}}
\newcommand{\Hh}{\mathcal{H}}
\newcommand{\Rr}{\mathcal{R}}
\newcommand{\Dd}{\mathsf{D}}
\newcommand{\Ll}{\mathsf{L}}
\newcommand{\Ii}{\mathcal{I}}
\newcommand{\Tt}{\mathcal{T}}
\DeclareMathOperator*{\argmin}{arg\,min}
\DeclareMathOperator*{\argmax}{arg\,max}
\hypersetup{pdftitle={Model 1023：基于规则信念的学习与决策模型},pdfauthor={},bookmarksdepth=3}
\title{Model 1023\\[0.5em]
\Large 学习与决策闭环下的认知隐状态生成模型\\[0.5em]
\large 规则信念、反馈调节与知识继承}
\author{}
\date{2026年10月24日}
\begin{document}
\maketitle
\thispagestyle{empty}


Model 1023描述学习者如何根据反馈调整对规则的信念、探索新的规则，并据此作出选择。
模型将当前参与判断和学习的规则称为活跃规则，以归一化规则信念表示它们在当前集合中的
相对可信度。各规则的信念之和为1；同一分布用于混合类别预测、选择局部搜索的起点，
以及确定容量已满时优先替换哪条规则。

反馈通过两项机制调节规则信念。正确反馈使已有信念比例在本次学习中完整保留，
错误反馈按个体保持率调节已有信念的集中程度，随后结合反馈似然更新。
反馈似然是各规则预测当前实际反馈的概率，用来评价它们对这次经验的解释程度。
新规则进入已满的活跃集合时，正确反馈后的新规则接替被移除规则的信念份额；
错误反馈后的新规则获得一个均匀试用份额。部分正确反馈取两者之间的值。
这两项机制分别决定已有判断如何调整，以及新判断获得多少后续检验机会。

模型由决策与学习两个模块构成。决策模块根据当前状态和刺激生成选择；学习模块根据选择与
反馈更新状态。认知状态包含活跃规则及其信念、判断精度、控制搜索的失败累积量，
以及条件3所需的类别配对信念。新规则继承当前规则的信念加权平均精度；条件3还继承
综合配对知识。条件1、2各有8项自由参数，条件3有9项。

本文先定义任务、规则空间与认知状态，再说明逐试次的决策和学习过程。
随后构造单试次生成核、轨迹推断算子和状态演化算子，说明个体参数的估计方法，
以及行为和口述分别检验模型的哪些内容。

\paragraph{阅读提示}
第\ref{sec:overview}节概括机制分工，第\ref{sec:mechanisms}节说明逐试次机制。
规则信念的保持与更新见第\ref{sec:retention}、\ref{sec:belief-update}节，
新规则的信念分配见第\ref{sec:admission}节，精度与配对继承见第\ref{sec:transfer}节。
参数定义与选参方法见第\ref{sec:inference}节，过程预测与检验范围见第\ref{sec:validation}节。

\paragraph{数据用途}
个体参数依据训练段选择似然和训练段自主生成表现确定，后段行为用于冻结后的评价。
自主生成通过模拟选择获得相应环境反馈，并形成完整学习路径。
口述仅用于检验固定参数下的内部表征，评价保留相容、冲突、未确定质量和目录覆盖率。
部分行为与口述资料已参与机制开发，相关结果按开发检验解释。

\clearpage
{\small\setlength{\parskip}{0pt}\tableofcontents}
\clearpage

\section{框架总览}
\label{sec:overview}

\subsection{两个心理模块与两个数学算子}
模型将一次学习试次分成决策与学习两个阶段。决策模块接收当前状态和刺激，产生选择；
学习模块结合该选择及反馈，更新下一题使用的状态：
\begin{equation}
(S_t,\bm x_t)\xrightarrow{\Dd_\theta}(m_t,y_t),\qquad
(S_t,m_t,y_t,r_t)\xrightarrow{\Ll_\theta}S_{t+1}.
\label{eq:loop}
\end{equation}
$S_t$是第$t$题作答前的认知状态，$\bm x_t$是物理刺激，$y_t$是选择，$r_t$是反馈。
$m_t$保存本题的内部知觉和规则预测，使学习使用与作答一致的信息。
$\theta$为同一被试在整段任务中保持固定的参数。

研究者在这个闭环上定义两个数学算子。\textbf{轨迹推断算子}$\Ii$结合实际行为，
更新各条可能内部路径的概率；\textbf{状态演化算子}$\Tt$对可能的选择和反馈求和，
给出模型自主学习的状态分布。粒子滤波是计算这些概率的数值方法，粒子代表研究者考虑的不同潜在路径。
认知容量$K$限制每条路径内的活跃规则数。

\subsection{参数、状态与试次记录}
参数规定学习过程的稳定属性，如搜索倾向、规则信念保持率和精度学习率。
状态随经验变化，记录当前考虑哪些规则、各规则的相对可信度，以及如何判断刺激。
本题记录$m_t$连接决策和学习，在反馈更新后完成其作用。

心理含义有两个层次：参数的\emph{操作定义}由公式给出；参数能否可靠地描述某个被试，
取决于数据量、估计精度、参数补偿关系和恢复检验。
例如，失败敏感度为零时，搜索概率固定为基础搜索概率，失败累积量的保留参数在此边界上失去行为作用。
第\ref{sec:parameters}节将这些边界与参数定义放在一起。

\subsection{各认知机制的分工}
决策模块依次完成知觉与选择，学习模块完成反馈评价、精度学习与规则管理。
条件3的反馈评价还包含对类别配对关系的学习。各机制通过当前状态和本题内部记录衔接。

{\small\noindent
\begin{tabularx}{\linewidth}{L{2.5cm}L{3.3cm}Y}
\toprule
所属阶段 & 机制 & 输入、输出和作用\\
\midrule
决策 & 知觉 & 根据物理刺激和个体知觉误差形成内部刺激。\\
决策 & 选择 & 按规则信念混合各规则的类别预测，生成选择概率。\\
学习 & 反馈评价 & 计算反馈似然，调节并更新规则信念；条件3同时更新配对信念。\\
学习 & 精度学习 & 根据对数反馈似然的梯度调整各活跃规则的判断精度。\\
学习 & 规则管理 & 更新失败累积量，搜索新规则，完成容量替换、入场信念分配与知识继承。\\
\bottomrule
\end{tabularx}}


\subsection{规则信念如何连接选择、搜索与保留}
规则信念表示学习者在当前活跃规则之间如何分配可信度。信念较高的规则对作答影响更大，
也更可能成为局部搜索的修改起点。新规则进入已满的活跃集合时，信念最低的已有规则被替换。
因此，规则发现、保留和使用通过同一个信念分布联系起来。

每次反馈首先调节已有信念的保持程度，再通过各规则的反馈似然改变它们的相对比例。
这一过程表示学习者根据判断是否成功，调整对已有经验的依赖。规则集合改变时，模型为
新规则分配试用份额，幸存规则按原比例分配剩余概率质量，使新集合的信念仍然归一化。

失败累积量决定搜索何时发生，判断精度决定各规则如何把类别距离转化为预测概率。
规则信念、搜索控制和精度学习共同决定后续选择。精度与配对继承把已经形成的任务知识
带入新规则的判断过程。

\section{实验设计、规则空间与认知状态}
\label{sec:state}

\subsection{实验条件与观测}
第$t$个有效试次包含四维归一化长度、实际选择和实际反馈：
\begin{equation}
\bm x_t\in[0,1]^4,\qquad y_t\in\{1,\ldots,C\},\qquad r_t\in\Rr.
\end{equation}
四维对应实验中的四个部位，在整段任务中保持相同顺序。条件1为二分类二值反馈，
$C=2$、$\Rr=\{0,1\}$；条件2为四分类二值反馈，$C=4$、$\Rr=\{0,1\}$；
条件3为四分类三值反馈，$C=4$、$\Rr=\{0,0.5,1\}$。
条件3选对种类得1分，同科但种类错误得0.5分，另一科得0分。

任务输入$u_t$包含刺激、任务条件及环境反馈分布$e_{u_t}(r\mid\bm x_t,y)$。
真实类别与按键映射用于环境评分及数据坐标转换；学习者根据自己的选择和获得的反馈学习。
条件3的真实类别、统一反应编号和实际按键分别记录，未知配对从均匀信念开始。
漏答试次跳过，模型按有效试次的原始顺序递推，并保留其原始试次编号。

\subsection{规则目录与活跃集合}
规则目录为$\Hh=\{1,\ldots,J\}$，二分类共有29条，四分类共有152条，分组见附录\ref{app:catalogue}。
每条规则规定一组判断、类别区域及固定类别标签。当前活跃集合$\A_t$由容量参数$K$约束：
\begin{equation}
\mu(h)=1/J,\qquad \A_t\subset\Hh,\qquad 1\leq|\A_t|\leq K<J.
\label{eq:capacity}
\end{equation}
$\mu$为基础规则权重，当前估参允许$K\in\{2,3,4,6,8\}$。
模型从目录中发现规则，目录的几何形式与标签共同限定其可表达策略。


\subsection{决策就绪状态与初始化}
把前题反馈引起的学习和规则管理完成后，本题作答前的状态写成
\begin{equation}
S_t=\left(\A_t,\{\pi_t(h),\beta_t(h)\}_{h\in\A_t},F_t,
\{q_t(k\mid h)\}_{h\in\A_t,\ k=1,2,3}\right).
\label{eq:state}
\end{equation}
$\pi_t(h)$为规则信念，$\beta_t(h)$为判断精度，$F_t$为失败累积量。
规则信念满足$\pi_t(h)\geq0$、$\sum_{h\in\A_t}\pi_t(h)=1$。
$q_t(k\mid h)$为给定规则$h$的条件配对信念，仅条件3使用，且每条规则的三种配对概率之和为1。
条件3也可用联合分布$\Pi_t(h,k)=\pi_t(h)q_t(k\mid h)$表示这两层信念。

首题抽取$h_1\sim\mu$，设置$\A_1=\{h_1\}$、$\pi_1(h_1)=1$、
$\beta_1(h_1)=\beta_0$、$F_1=0$；条件3令$q_1(k\mid h_1)=1/3$。
首题完成后才根据其反馈启动搜索控制。

本题内部记录为
\begin{equation}
m_t=\left(\widetilde{\bm x}_t,\{d_{t,h,c}\},\{p_{t,h}(c)\},\{\pi_t(h)\}\right),
\end{equation}
分别保存内部刺激、类别距离、规则类别预测和作答前信念。
学习模块使用这份记录与本题反馈完成更新，形成下一有效试次的$S_{t+1}$。
$S_t$与$S_{t+1}$对应相同的时点：前一试次学习已完成，当前试次尚未作答。

\section{认知机制设计}
\label{sec:mechanisms}

\subsection{决策模块：从当前状态生成选择}
\label{sec:decision}

\subsubsection{知觉：从物理刺激形成内部刺激}
决策模块的输入是$S_t$与$\bm x_t$，输出本题内部记录和选择，其概率规律记为
$\Dd_\theta(dm,y\mid S,\bm x)$。它先形成内部刺激，再计算各规则预测，并混合生成选择。
各规则共享同一次知觉结果：
\begin{equation}
\widetilde x_{t,j}=\clip_{[0,1]}(x_{t,j}+\epsilon_{t,j}),\qquad j=1,\ldots,4.
\label{eq:perception}
\end{equation}
$\clip$将越界值限制在指定区间，$\epsilon_{t,j}$表示第$j$个特征的知觉误差。
误差采用正态或均匀分布，其分布形式与参数由独立知觉测量确定，并与个体和特征对应。
行为参数估计期间，这些测量参数保持固定。

\subsubsection{类别预测：从区域距离计算类别概率}
规则$h$将刺激空间划分为类别区域$\mathcal R_{h,c}$。
模型用内部刺激到各区域的最短欧氏距离计算类别概率：
\begin{align}
d_{t,h,c}&=\inf_{\bm z\in\mathcal R_{h,c}}\norm{\widetilde{\bm x}_t-\bm z}_2,\\
p_{t,h}(c)&=\frac{\exp[-\beta_t(h)d_{t,h,c}]}
{\sum_{a=1}^C\exp[-\beta_t(h)d_{t,h,a}]}.
\label{eq:category}
\end{align}
刺激位于某区域内时距离为0；类别有多个分离分量时，取最近分量的距离。
$\beta$越大，概率越集中在距离较小的类别。知觉误差改变内部刺激，
精度改变类别距离对判断的影响，两者作用于不同环节。


\subsubsection{混合选择：按规则信念整合类别预测}
作答前信念$\pi_t(h)$直接规定各规则对本题选择的相对贡献。
将各规则的类别预测按信念加权相加，得到选择概率：
\begin{equation}
p_t(c\mid S_t,\widetilde{\bm x}_t)=\sum_{h\in\A_t}\pi_t(h)p_{t,h}(c).
\label{eq:choice}
\end{equation}
规则信念和各规则的类别概率分别归一化，因此混合后的类别概率之和也为1。
例如，两条规则的信念为0.75和0.25，分别给类别1以0.8和0.2的概率，
混合后选择类别1的概率为$0.75\times0.8+0.25\times0.2=0.65$。

自主生成从该分布抽取$y_t$；行为计分读取实际选择的概率。
信念分散时，多条规则共同影响回答；信念集中时，一条规则主导回答。
额外随机作答率固定为0，行为变异来自知觉、规则搜索和概率选择。

有效参与量$N_{\mathrm{eff},t}=1/\sum_h\pi_t(h)^2$描述信念的集中程度，
其值位于1与$|\A_t|$之间，可为非整数。各规则信念相等时，它等于活跃规则数；
一条规则占据绝大部分信念时，它接近1。活跃规则数与有效参与量分别描述容量占用和概率分配。

\subsection{学习模块：从反馈形成下一状态}
\label{sec:learning}

\subsubsection{学习顺序：用本题预测评价反馈}
学习模块接收$S_t,m_t,y_t,r_t,u_t$，其输出分布记为
$\Ll_\theta(dS'\mid S,m,y,r,u)$。
它用作答时的类别预测计算反馈似然，再更新规则信念、精度与失败累积量，
随后提出候选、处理容量并初始化新规则。条件3同时更新配对知识。

\textbf{反馈评价使用作答时的预测。}信念更新与精度梯度共享$m_t$中的知觉、距离和类别概率。
精度学习后的值用于后续判断和新规则继承。
因此，反馈先评价本题作答时的预测，再改变后续试次的判断。

\subsubsection{二值任务的反馈似然}
反馈似然$\ell_t(h)$衡量规则$h$对当前实际反馈的预测程度。
正确反馈对应所选类别，错误反馈对应其余类别的总概率：
\begin{equation}
\ell_t(h)=
\begin{cases}
p_{t,h}(y_t),&r_t=1,\\
1-p_{t,h}(y_t),&r_t=0.
\end{cases}
\label{eq:binary-feedback}
\end{equation}
因此，四分类中的错误反馈评价三个其他类别的总概率。
所有活跃规则都接收这次反馈评价，各规则的反馈似然由其自己的预测决定。

\subsubsection{规则信念保持：根据反馈调整已有判断}
\label{sec:retention}
模型假设，成功的判断使学习者更充分地保持已有信念，失败则降低对已有判断的依赖。
设$\delta\in[0,1]$为个体的基础规则信念保持率，本题有效保持率为
\begin{equation}
\delta_t=\delta+(1-\delta)r_t.
\label{eq:effective-retention}
\end{equation}
正确反馈$r_t=1$时，$\delta_t=1$；错误反馈$r_t=0$时，$\delta_t=\delta$；
半分反馈$r_t=0.5$时，$\delta_t=(1+\delta)/2$。
例如$\delta=0.8$时，三种反馈对应的有效保持率依次为1、0.8和0.9。

在结合本题反馈似然之前，先对已有信念作幂变换并归一化：
\begin{equation}
\bar\pi_t(h)=\frac{\pi_t(h)^{\delta_t}}
{\sum_{v\in\A_t}\pi_t(v)^{\delta_t}}.
\label{eq:belief-retention}
\end{equation}
当$\delta_t=1$时保留原有比例；当$0<\delta_t<1$时，正信念之间的比值更接近1。
当$\delta_t=0$时，约定全部活跃规则的$\bar\pi_t$均为$1/|\A_t|$，包括此前信念为零的规则。
这一步完成后，再由反馈似然决定哪些规则应获得更高的更新后信念。

$\bar\pi_t$是收到反馈后用于学习的中间分布。本题选择已由$\pi_t$生成，
当前反馈通过$\delta_t$和后续似然更新共同改变下一题的状态。
这是一项反馈依赖的信念学习假设，其心理解释需要结合行为与内部状态检验。

\subsubsection{条件3的配对信念与反馈似然}
\label{sec:pairing}
条件3的四个反应有三种可能的两两配对：
\begin{equation}
G_1=\{\{1,2\},\{3,4\}\},\quad G_2=\{\{1,3\},\{2,4\}\},\quad
G_3=\{\{1,4\},\{2,3\}\}.
\end{equation}
令$\operatorname{mate}_k(y)$为配对$G_k$中与$y$同科的另一个反应。
对于规则$h$和配对$k$，本题反馈似然为
\begin{equation}
\ell_t(h,k)=\begin{cases}
p_{t,h}(y_t),&r_t=1,\\
p_{t,h}(\operatorname{mate}_k(y_t)),&r_t=0.5,\\
1-p_{t,h}(y_t)-p_{t,h}(\operatorname{mate}_k(y_t)),&r_t=0.
\end{cases}
\label{eq:pair-likelihood}
\end{equation}
满分、半分和零分分别对应所选种类、同科另一种类和另一科的两个种类。


条件配对信念由保持率$\gamma$调整：
\begin{equation}
\bar q_t(k\mid h)=\frac{q_t(k\mid h)^\gamma}
{\sum_{j=1}^3q_t(j\mid h)^\gamma}.
\label{eq:pair-forgetting}
\end{equation}
$\gamma=0$时，每条规则的三种配对设为均匀分布；$\gamma=1$时完整保留原比例。
规则层使用式\eqref{eq:belief-retention}的$\bar\pi_t$，条件配对层使用$\bar q_t$，
两层分别归一化。规则对实际反馈的预测概率为
\begin{equation}
\ell_t(h)=\sum_k\bar q_t(k\mid h)\ell_t(h,k).
\label{eq:marginal-likelihood}
\end{equation}
模型据此形成规则与配对的联合更新：
\begin{align}
\Pi_t^+(h,k)&=
\frac{\bar\pi_t(h)\bar q_t(k\mid h)\ell_t(h,k)}
{\sum_{v\in\A_t}\sum_j\bar\pi_t(v)\bar q_t(j\mid v)\ell_t(v,j)},\\
\pi_t^+(h)&=\sum_k\Pi_t^+(h,k),\qquad
q_t^+(k\mid h)=\frac{\Pi_t^+(h,k)}{\pi_t^+(h)},\qquad
Q_t^+(k)=\sum_h\Pi_t^+(h,k).
\label{eq:pair-update}
\end{align}
这里$\Pi_t^+$是反馈后的联合信念，$\pi_t^+$是规则边际，$q_t^+$是各规则的条件配对信念。
规则边际为零时，约定其条件配对分布为均匀分布。
$Q_t^+$汇总全部已有规则的配对知识，供新规则继承。
精度学习使用的配对权重为$\bar q_t$，与本题反馈似然的计算时点一致。

两项保持率的分工是：$\delta_t$调节规则之间的信念比例，$\gamma$调节各规则内部的配对比例。
两层再通过同一反馈似然形成联合信念。条件1、2使用式\eqref{eq:binary-feedback}，
配对保持率$\gamma$仅在条件3估计。

\subsubsection{规则信念更新：结合已有判断与反馈似然}
\label{sec:belief-update}
三种条件均用本题反馈似然重新分配规则信念：
\begin{equation}
\pi_t^+(h)=\frac{\bar\pi_t(h)\ell_t(h)}
{\sum_{v\in\A_t}\bar\pi_t(v)\ell_t(v)}
=\frac{\pi_t(h)^{\delta_t}\ell_t(h)}
{\sum_{v\in\A_t}\pi_t(v)^{\delta_t}\ell_t(v)}.
\label{eq:belief-update}
\end{equation}
指数为零时按第\ref{sec:retention}节的均匀约定处理。
条件3的式\eqref{eq:belief-update}正是式\eqref{eq:pair-update}的规则边际，
因此每次反馈只形成一次规则信念更新。

在保持后信念相同时，反馈似然越大的规则越能解释实际反馈，获得的更新后信念也越大。
例如，保持后的两条规则信念为0.6和0.4，反馈似然分别为0.8和0.2，
则未归一化质量为0.48和0.08，更新后信念约为0.857和0.143。
两条规则均使用本次反馈；更新结果同时取决于已有比例和预测反馈的能力。

对于信念与似然均为正的两条规则$h,v$，更新还可写成对数比值的形式：
\begin{equation}
\log\frac{\pi_t^+(h)}{\pi_t^+(v)}
=\delta_t\log\frac{\pi_t(h)}{\pi_t(v)}
+\log\frac{\ell_t(h)}{\ell_t(v)}.
\label{eq:belief-odds}
\end{equation}
第一项保留此前的相对判断，第二项加入当前反馈对两条规则的区分。
此前经验通过当前信念持续影响学习，其影响随各题的有效保持率及规则集合变化而调整。
更新后的$\pi_t^+$用于搜索、容量替换和继承，随后经新规则入场分配形成$\pi_{t+1}$。

\subsubsection{精度学习：按反馈调整类别区分程度}
\label{sec:precision}
精度学习使用对数反馈似然对$\beta$的梯度，表示小幅改变精度时，对数反馈似然如何变化。
令$a_{t,h,c}$为类别$c$与实际反馈相容的程度。二值任务中，
$r=1$时$a_{t,h,c}=\ind[c=y_t]$，$r=0$时$a_{t,h,c}=\ind[c\ne y_t]$。
条件3对三种配对加权：
\begin{equation}
a_{t,h,c}=\sum_k\bar q_t(k\mid h)\ind[c\in Y_k(y_t,r_t)],
\end{equation}
其中，$\ind$为指示函数，条件成立时取1，否则取0；$Y_k$是式\eqref{eq:pair-likelihood}中与实际反馈对应的类别集合。
将本题类别预测按反馈相容程度加权并归一化，得到反馈条件下的类别分布：
\begin{equation}
p_{t,h}^{\mathrm{fb}}(c)=\frac{a_{t,h,c}p_{t,h}(c)}{\ell_t(h)}.
\end{equation}
固定本题知觉、距离和反馈前配对分布，梯度可写成两个平均距离之差：
\begin{equation}
g_t^\beta(h)=\left.\frac{\partial\log\ell_t(h;\beta)}{\partial\beta}\right|_{\beta_t(h)}
=\sum_c p_{t,h}(c)d_{t,h,c}-\sum_c p_{t,h}^{\mathrm{fb}}(c)d_{t,h,c}.
\label{eq:gradient}
\end{equation}
反馈相容类别的平均距离较小时，提高精度；较大时，降低精度。
更新方向由反馈与规则的实际预测共同确定。

判断精度按梯度作加法更新，并限制在$[0.1,25]$内：
\begin{equation}
\beta_t^+(h)=\clip_{[0.1,25]}\left[\beta_t(h)+10\eta g_t^\beta(h)\right].
\label{eq:beta}
\end{equation}
固定系数10规定学习率的数值尺度，$\eta$为个体参数，有效步长为$10\eta$。
每条活跃规则均按自己的梯度更新。有限步长和截断共同决定实际变化，
单步似然的实际增减需由该更新结果计算。

\subsubsection{失败累积量与搜索触发}
失败累积量$F_t$控制下一题前是否探索新规则。完全成功使已有失败量衰减，
不完全成功则按失分程度增加：
\begin{equation}
F_{t+1}=\begin{cases}\lambda F_t,&r_t=1,\\F_t+(1-r_t),&r_t<1.\end{cases}
\label{eq:failure}
\end{equation}
因此条件3的半分增加0.5，零分增加1。$\lambda\in[0,1]$表示一次完全成功之后保留多少失败量。
搜索概率为
\begin{equation}
s_{t+1}=\sigmoid\!\left[\logit(p_0)+a_eF_{t+1}\right],\qquad
\sigmoid(z)=\frac1{1+e^{-z}},\quad\logit(p)=\log\frac{p}{1-p}.
\label{eq:search}
\end{equation}
$p_0\in(0,1)$为失败量为零时的基础搜索概率，$a_e\geq0$为失败敏感度。
每次抽取一个概率为$s_{t+1}$的搜索事件，发生时提出一条新规则。
完全成功后的基础搜索机会允许模型继续探索。

\subsubsection{局部结构搜索与全局搜索}
\label{sec:proposal}
候选范围为$\mathcal U_t=\Hh\setminus\A_t$。
搜索发生后，以固定个体比例$g_0$采用全局路线，从$\mu$在$\mathcal U_t$上的归一化分布提出候选。

其余情况采用局部路线，按反馈更新后的规则信念抽取父规则$h_{\mathrm{parent}}$：
\begin{equation}
\Pr(h_{\mathrm{parent}}=h\mid\text{局部搜索},S_t,m_t,y_t,r_t)=\pi_t^+(h).
\label{eq:parent}
\end{equation}
信念较高的规则更可能成为进一步修改的起点。

结构邻域$\mathcal N_t(h)\subset\mathcal U_t$包含对规则$h$作一次合法结构修改后得到的规则。
编辑距离$d_{\mathrm{edit}}(h,u)$记录到候选$u$的最小单步成本；更换特征、
改变比较关系及标准参照、修改参照值、改写条件树一条下级判断的成本依次为1、1.5、1、2。
编辑保留其余判断和相应的类别输出布局，且结果必须在固定目录内。
非空邻域中的提议概率为
\begin{equation}
Q_{\mathrm{edit},t}(u\mid h)=
\frac{\ind[u\in\mathcal N_t(h)]e^{-d_{\mathrm{edit}}(h,u)}}
{\sum_{v\in\mathcal N_t(h)}e^{-d_{\mathrm{edit}}(h,v)}}.
\end{equation}
成本越低的修改越容易提出。

当结构邻域没有非活跃规则时，局部搜索依据固定类别标签下的功能相似性提出候选。
功能距离$d_{\mathrm{func}}(h,u)$是两条规则在均匀参考刺激上给出不同类别标签的比例。
该距离以100000个共同参考刺激估计，随机种子为0。功能搜索的提议分布为
\begin{equation}
Q_{\mathrm{func},t}(u\mid h)=
\frac{\ind[u\in\mathcal U_t]\mu(u)e^{-d_{\mathrm{func}}(h,u)/0.10}}
{\sum_{v\in\mathcal U_t}\mu(v)e^{-d_{\mathrm{func}}(h,v)/0.10}}.
\end{equation}
记$Q_{\mathrm{local},t}$为有邻居时的结构分布、无邻居时的功能回退分布，
则搜索事件内的完整候选分布为
\begin{equation}
Q_t(u)=g_0\frac{\mu(u)}{\sum_{v\in\mathcal U_t}\mu(v)}
+(1-g_0)\sum_{h\in\A_t}\pi_t^+(h)Q_{\mathrm{local},t}(u\mid h),
\quad u\in\mathcal U_t.
\label{eq:proposal}
\end{equation}
新候选的概率由当前活跃状态和固定规则关系确定。
局部搜索在邻域非空时按结构编辑距离抽样，邻域为空时按功能距离抽样。

\subsubsection{容量替换与规则保留}
\label{sec:replacement}
没有搜索事件时，$\A_{t+1}=\A_t$。
发生搜索并提出$u$后，有空位就加入；容量已满就从已有规则中移除信念最低的一条：
\begin{equation}
\A_{t+1}=\begin{cases}
\A_t\cup\{u\},&|\A_t|<K,\\
(\A_t\setminus\{h^-\})\cup\{u\},&|\A_t|=K,
\end{cases}
\qquad h^-\in\argmin_{h\in\A_t}\pi_t^+(h).
\label{eq:replacement}
\end{equation}
多条规则的信念并列最低时，等概率抽取其中一条。
比较对象为已有规则，新候选在搜索发生时获得一次入集机会。

因此，规则保留取决于当前集合是否出现新候选，以及它在已有规则中的相对信念。
信念较低的规则可继续留在尚未搜索的集合中，其作答贡献由当前较低的信念决定。
规则集合从首题的一条逐渐增加，到达容量后通过替换改变组成。


\subsubsection{新规则入场：按反馈分配初始信念}
\label{sec:admission}
新规则需要获得初始信念，才能参与下一题判断。模型根据集合是否已满，以及刚完成试次的
反馈，为新规则分配概率质量。没有搜索事件时，直接令$\pi_{t+1}=\pi_t^+$。

发生搜索并提出新规则$u$时，记其入场份额为$\nu_{t+1}$。
当活跃集合尚未达到容量时，设$n_t=|\A_t|$，新规则获得新集合的均匀份额$1/(n_t+1)$。
当集合已满时，设$h^-$为本次被移除的规则，其份额按反馈调节：
\begin{equation}
\nu_{t+1}=\begin{cases}
\displaystyle\frac{1}{n_t+1},&n_t<K,\\[0.5em]
\displaystyle r_t\pi_t^+(h^-)+\frac{1-r_t}{K},&n_t=K.
\end{cases}
\label{eq:entry-share}
\end{equation}
满容量时，答对后的新规则接替被移除规则的份额；答错后获得$1/K$的试用份额；
半分反馈取二者的平均。例如$K=4$且被移除规则的信念为0.001时，正确、半分和错误
反馈后的新规则份额依次为0.001、0.1255和0.25。该分配表示判断失败时更充分地尝试新解释。

记幸存规则集合为$\A_t^{\mathrm{keep}}=\A_t\cap\A_{t+1}$。
新规则获得$\nu_{t+1}$，幸存规则按更新后的相对比例分配剩余质量：
\begin{equation}
\pi_{t+1}(u)=\nu_{t+1},\qquad
\pi_{t+1}(h)=(1-\nu_{t+1})
\frac{\pi_t^+(h)}{\sum_{v\in\A_t^{\mathrm{keep}}}\pi_t^+(v)},
\quad h\in\A_t^{\mathrm{keep}}.
\label{eq:entry-belief}
\end{equation}
这一步保持幸存规则之间的信念比值，并使新集合内的总信念为1。
满容量且答对时，幸存规则的绝对信念也保持原值。新规则从下一题开始接受自身的反馈检验。
式\eqref{eq:entry-share}中的$r_t$始终是新规则入场之前刚完成试次的反馈。

被移除规则的信念至多为$1/K$，因此错误后的均匀试用份额至少与其原份额相同。
若被移除规则的信念恰为零，答对后的新规则也得到零信念；有效保持率大于零时，
后续乘法更新保留该零值。有效保持率为零时，活跃规则重新获得均匀的保持后分布。
这些边界说明，规则在场与它对回答的实际贡献需要分别检查。

\subsubsection{精度与配对继承：延续已有任务知识}
\label{sec:transfer}
保留规则延续本题学习后的判断精度，新规则继承全部已有活跃规则的信念加权平均精度：
\begin{equation}
\bar\beta_t=\sum_{h\in\A_t}\pi_t^+(h)\beta_t^+(h),\qquad
\beta_{t+1}(h)=\begin{cases}\beta_t^+(h),&h\in\A_t^{\mathrm{keep}},\\
\bar\beta_t,&h=u.\end{cases}
\label{eq:transfer}
\end{equation}
参与平均的是\textbf{容量替换之前的全部活跃规则}，包括本次随后被移除的规则；
信念与精度均已吸收本题反馈。信念非负且和为1，因此继承值位于参与平均的最小精度与最大精度之间。
首题规则使用$\beta_0$。

这一继承假设任务中形成的刺激区分能力能够在规则之间延续。
例如，已有规则的更新后精度为6和10、信念为0.75和0.25时，新规则精度为7。
其后每条规则仍按自己的几何预测和反馈调整精度。

条件3中，保留规则延续自己的条件配对信念，新规则继承式\eqref{eq:pair-update}的整体配对边际：
\begin{equation}
q_{t+1}(k\mid h)=\begin{cases}q_t^+(k\mid h),&h\in\A_t^{\mathrm{keep}},\\
Q_t^+(k),&h=u,\end{cases}\qquad
\Pi_{t+1}(h,k)=\pi_{t+1}(h)q_{t+1}(k\mid h).
\label{eq:pair-transfer}
\end{equation}
$Q_t^+$同样在容量替换前计算，以更新后的规则信念汇总各规则的配对知识。
局部与全局提出的新规则采用相同的信念分配及继承方式。

规则离开活跃集合时，其专属信念、精度与条件配对状态随之结束。
将来重新提出时，按届时的反馈与活跃状态分配入场信念，并继承精度及配对知识。

\subsubsection{一次完整有效试次}
\begin{enumerate}
\item 从$S_t$读取活跃规则、信念、精度及配对信念，抽取本题内部刺激。
\item 计算各规则的类别预测，按$\pi_t$混合，生成选择或计算实际选择的概率。
\item 接收实际反馈；自主生成时，由环境对模拟选择评分。
\item 根据$r_t$确定有效保持率$\delta_t$；用作答时的预测计算反馈似然，更新规则信念。
条件3同时调整条件配对信念并完成一次联合更新。
\item 用同一份作答时预测更新判断精度，更新失败累积量，形成下一题前的搜索概率。
\item 搜索发生时按局部／全局路线提出新规则；容量已满时替换信念最低的已有规则。
\item 根据本题反馈为新规则分配入场信念，按更新后的信念继承平均精度；条件3继承综合配对知识。
\item 将新集合及其归一化信念、精度、失败累积量和配对信念组成$S_{t+1}$。
\end{enumerate}
反馈信念更新与精度学习使用原有活跃集合。新规则的初始状态由已有知识与入场分配确定，
其自身的首次反馈学习发生在下一有效试次。

\section{试次级隐状态生成}
\label{sec:kernel}

\subsection{联合生成核}
决策、环境反馈与学习共同规定从当前状态出发，下一状态及观测如何联合产生：
\begin{equation}
k_\theta(dS',y,r\mid S,u)=
\int \Dd_\theta(dm,y\mid S,\bm x)\,
e_u(r\mid\bm x,y)\,
\Ll_\theta(dS'\mid S,m,y,r,u).
\label{eq:kernel}
\end{equation}
这里的积分消去内部试次记录 $m$，并保留该记录在决策与学习之间的依赖关系。
核使用概率测度记法，以同时涵盖连续的知觉和精度、离散的规则集合及确定性的状态更新。
固定 $S,u$ 后，对 $S'$ 积分并对全部 $y,r$ 求和的结果为1。

\subsection{生成、计分与内部推断}
自主生成从 $\Dd$ 抽取选择，由环境产生相应反馈，再经 $\Ll$ 更新状态。
观测计分固定实际选择与反馈，对可能的知觉、规则集合和其他状态进行积分。
内部推断则以这些观测重新分配潜在路径的概率。三者共享同一组参数与更新机制。

当实验评分确定时，对有效记录有 $e_u(r_t\mid\bm x_t,y_t)=1$。
因此研究者的行为似然由选择概率构成，反馈作为已给定的学习输入改变下一状态。
规则的反馈似然 $\ell_t(h)$ 用于更新规则信念、判断精度和配对信念；粒子路径的观测权重按外显选择更新。
这一区分保证同一反馈的认知更新与行为序列计分各有明确位置。

\section{学习轨迹推断与状态分布演化}
\label{sec:operators}

\subsection{轨迹推断算子}
令 $b$ 为当前决策就绪状态的概率分布，$A'$ 为下一状态空间中的可测集合。
给定实际观测 $o=(y,r)$，定义未归一化算子
\begin{equation}
(\Ii_{\theta,o}^{u}b)(A')=\int k_\theta(A',o\mid S,u)\,b(dS).
\label{eq:inference}
\end{equation}
它按实际观测的生成概率重加权路径，并把路径推进到学习后的状态。
记 $\langle1,\zeta\rangle$ 为测度 $\zeta$ 的总质量，则
\begin{equation}
P_\theta(o\mid u,b)=\langle1,\Ii_{\theta,o}^{u}b\rangle,\qquad
b_{t+1}=\frac{\Ii_{\theta,o_t}^{u_t}b_t}
{\langle1,\Ii_{\theta,o_t}^{u_t}b_t\rangle}.
\label{eq:filter}
\end{equation}
第一式给出观测概率，第二式给出观察本题并完成学习后的状态后验。
未归一化的 $\Ii$ 对输入测度是正线性的；除以总质量的归一化过滤映射一般为非线性。

对给定输入序列，完整观测序列概率可写为
\begin{equation}
P_\theta(o_{1:T}\mid u_{1:T})=
\left\langle1,\Ii_{\theta,o_T}^{u_T}\cdots\Ii_{\theta,o_1}^{u_1}b_1\right\rangle.
\label{eq:sequence-operator}
\end{equation}
算子从右向左按时间作用。式\eqref{eq:filter}给出各时点的过滤边际；
完整路径的联合后验还需保留路径间的祖先关系。
利用未来观测更新过去状态属于平滑推断。逐试次状态分析采用在线过滤，并明确区分选择前与选择后时点。
输入若由历史自适应决定，则按已知实验策略条件化，记录其与此前观测的依赖。

\subsection{状态演化算子}
尚未给定实际选择和反馈时，对所有可能观测求和得到
\begin{equation}
\Tt_{\theta,u}=\sum_y\sum_{r\in\Rr}\Ii_{\theta,(y,r)}^{u},\qquad
\langle1,\Tt_{\theta,u}b\rangle=1.
\label{eq:evolution}
\end{equation}
状态演化算子把当前状态分布映射为完成一次决策与学习后的预测分布。
给定预先确定的刺激和反馈协议序列，
\begin{equation}
b_{T+1}^{\mathrm{gen}}=\Tt_{\theta,u_T}\cdots\Tt_{\theta,u_1}b_1.
\end{equation}
自主轨迹模拟通过抽样近似这一传播。
模拟中的反馈必须由模拟回答和实际任务评分规则产生。
自适应输入协议则把输入策略一并放入前向生成过程。

\subsection{各层的用途与分析条件}
学习核 $\Ll$ 描述给定一次内部记录、选择和反馈后，单个认知状态怎样更新。
$\Ii$ 描述观察实际行为后，研究者的状态分布怎样改变；$\Tt$ 描述尚未观察行为时，
所有可能的决策和学习结果形成怎样的状态分布。这些对象由同一单试次核连接。

在同一状态空间和兼容输入协议下，比较 $\Tt_{u_b}\Tt_{u_a}$ 与 $\Tt_{u_a}\Tt_{u_b}$
可定义信息顺序对学习分布的影响。稳态与谱分析还要求明确平稳输入策略、状态空间及有限维近似。
认知状态包含连续的规则信念、判断精度和失败累积量；长期分析需规定这些变量的取值处理及近似误差。
长期谱、稳态及前瞻干预的分析结果尚待计算。

\section{个体参数、数值推断与估计目标}
\label{sec:inference}

\subsection{八项共同参数与条件3的配对参数}
\label{sec:parameters}
条件1、2的个体参数为
\begin{equation}
\theta_s=(p_{0,s},a_{e,s},\lambda_s,g_{0,s},K_s,\delta_s,\beta_{0,s},\eta_s).
\label{eq:parameters}
\end{equation}
条件3再加入$\gamma_s$。以下省略被试下标；各参数在一次拟合的整段试次内固定，
其共同作用产生随时间变化的状态和行为。

{\small
\begin{longtable}{L{1.3cm}L{3.0cm}L{5.2cm}L{4.7cm}}
\toprule
参数 & 名称与范围 & 在模型中的作用 & 解释时的边界\\
\midrule\endhead
$p_0$ & 基础搜索概率，$(0,1)$ & 失败累积量为零时的搜索概率。 & 实际搜索率还由$a_e$与$F_t$决定。\\
$a_e$ & 失败敏感度，$[0,\infty)$ & 失败量对搜索对数优势的增益。 & 为零时，$\lambda$失去行为作用。\\
$\lambda$ & 失败累积量保留，$[0,1]$ & 完全成功后保留的失败量比例。 & 作用依赖成功与失败的时序。\\
$g_0$ & 全局搜索比例，$[0,1]$ & 搜索发生时选择全局目录的概率。 & 局部路线存在邻域为空时的功能回退。\\
$K$ & 活跃容量，整数$\geq2$ & 同时保留的规则数上限。 & 是当前规则表示下的有效容量。\\
$\delta$ & 规则信念保持，$[0,1]$ & 零分反馈时的信念保持率；本题有效值为$\delta+(1-\delta)r_t$。 & 正确反馈时有效值为1；其作用依赖反馈时序。\\
$\beta_0$ & 初始精度，$[0.1,25]$ & 首题规则的类别距离敏感程度。 & 后续精度还受学习和继承影响。\\
$\eta$ & 精度学习率，$[0,1]$ & 控制$10\eta$倍反馈梯度的更新幅度。 & 为零时，全部规则保持$\beta_0$。\\
$\gamma$ & 配对信念保持，$[0,1]$ & 条件3中各规则的条件配对分布保持率。 & 仅条件3估计；与精度存在补偿可能。\\
\bottomrule
\end{longtable}}

上述含义界定参数在模型中的作用。把点估计进一步解释为稳定个体特质，需要检查：
不同起点能否找到相近的高似然区域、相近分数下参数是否分散、模拟真值能否恢复，
以及关键内部状态是否对参数变化稳健。已知零作用参数在对应参数点单独标记。

当$\eta=0$时，反馈引起的精度增量为零。首题规则的精度为$\beta_0$，
后续新规则继承的平均精度也为$\beta_0$，因此全部活跃规则在整段任务中保持相同精度。
在此边界上，规则信念和活跃集合仍可随反馈变化。

\subsection{粒子滤波：积分可能的内部路径}
第$n$个粒子表示一条潜在路径，作答前权重为$W_t^{(n),\mathrm{pre}}$。
由其状态和抽取的内部刺激得到选择概率$p_t^{(n)}(c)$，研究者的作答前预测为
\begin{equation}
P_t(c)=\sum_{n=1}^{N_{\mathrm{PF}}}W_t^{(n),\mathrm{pre}}p_t^{(n)}(c),\qquad
\sum_nW_t^{(n),\mathrm{pre}}=1.
\label{eq:pf-predict}
\end{equation}
实际选择出现后，路径权重更新为
\begin{equation}
W_t^{(n),\mathrm{choice}}=
\frac{W_t^{(n),\mathrm{pre}}p_t^{(n)}(y_t)}{P_t(y_t)}.
\label{eq:pf-update}
\end{equation}
更能预测实际回答的路径获得更高后验权重；所有路径再根据实际反馈完成各自的学习。
有效样本量（effective sample size，ESS）衡量粒子权重的集中程度。
当$\mathrm{ESS}=1/\sum_n(W_t^{(n),\mathrm{choice}})^2$低于预先设定的阈值时进行重采样，
子粒子继承认知状态，并产生独立的后续随机分支。
粒子数、阈值和重复次数属于数值设置，按具体协议记录。

可在指定时点汇总规则在场概率、作答贡献和精度分布。
固定参数下的路径区间反映潜在过程的不确定性；独立计算重复的标准误反映数值误差；
个体参数的不确定性需另行传播。三者在图表中分别注明。

\subsection{训练似然：评价作答前选择预测}
训练目标为实际选择的平均负对数似然（negative log-likelihood，NLL）：
\begin{equation}
\mathrm{NLL}_{\mathrm{train}}(\theta_s)=
-\frac1{|\mathcal J_{\mathrm{train}}|}\sum_{t\in\mathcal J_{\mathrm{train}}}
\log P_t(y_t;\theta_s).
\label{eq:objective}
\end{equation}
$\mathcal J_{\mathrm{train}}$为预先固定的训练题集合，$P_t$在观察本题选择之前计算。
计分覆盖全部有效训练题，包括首题；本题选择和反馈在计分后进入学习。
多次独立滤波先逐题平均选择概率，再取对数得到参数比较分数。

在同一生成模型和任务协议下，各题条件选择概率的乘积等于行为序列的生成概率。
因而提高精确序列似然，就提高模型对观测序列的概率。
粒子滤波近似计算状态积分，实际比较仍受有限粒子数和参数搜索范围影响。
这一阶段确定已检查候选中的高似然参数点。其可靠性通过多起点搜索、独立随机种子的评分确认，
以及已知生成参数的模拟恢复检验来评价。参数冻结后，再使用独立数值重复评价行为预测。

\subsection{受训练似然约束的自主生成选参}
\label{sec:selection}
参数估计同时检查逐题选择预测与从任务起点自主学习的表现。训练选择似然先确定一组
能够解释实际行为的候选，训练自主生成再评价这些候选形成的学习路径分布。
全部选参步骤只使用预先固定的训练段。

自主路径的行为特征包括每32题的正确率，以及给定真实类别时选择各反应的频率。
设训练段共有$B$个时间块，类别条件反应频率共有$D$项，特征向量为
\begin{equation}
\phi(\bm y)=\left(\sqrt{\frac{0.5}{B}}a_1,\ldots,
\sqrt{\frac{0.5}{B}}a_B,\sqrt{\frac{0.5}{D}}b_1,\ldots,
\sqrt{\frac{0.5}{D}}b_D\right).
\label{eq:behavior-features}
\end{equation}
$a_j$为第$j$个时间块的正确率，末块不足32题时按实际题数计算；
$b_j$为某真实类别条件下选择某反应的比例。
刺激序列中出现的每个真实类别都保留全部$C$种反应频率，未出现的真实类别不构成该特征。
上述缩放使两类特征按维数平均后，在距离平方中分别使用0.5的权重。
真实类别用于环境评分及行为汇总，学习者接收其选择对应的反馈。

设$\bm f_m$为第$m$条自主路径的特征，$\bm f_{\mathrm{obs}}$为真实训练行为特征，
用Energy评分衡量路径分布与观察行为的接近程度：
\begin{equation}
\widehat{\mathcal E}(\theta)=
\frac1M\sum_{m=1}^M\norm{\bm f_m-\bm f_{\mathrm{obs}}}_2
-\frac{1}{2M(M-1)}\sum_{m\ne n}\norm{\bm f_m-\bm f_n}_2.
\label{eq:energy}
\end{equation}
第一项衡量模拟路径与观察行为的距离，第二项计入模拟分布自身的离散程度。
评分越小，生成分布与观察行为越相容。第二项排除路径与自身的比较，
按不同路径对构成U统计量。全部自主路径参与评分。

记$\theta_{\mathrm N}$为训练NLL最优候选，$\theta_{\mathrm G}$为生成表现候选。
在已确认的参数集合中，先保留训练NLL不超过最优值0.02的候选，再选择其中训练Energy最小者。
用于候选排序的每个参数点生成256条训练路径。

提议参数随后接受独立随机种子的两项复核。令$\Delta_{\mathrm{NLL}}$和
$\Delta_{\mathcal E}$分别为$\theta_{\mathrm G}$减去$\theta_{\mathrm N}$的训练评分差值，
仅当下列两项同时成立时接受提议：
\begin{equation}
\Delta_{\mathrm{NLL}}+2\operatorname{MCSE}(\Delta_{\mathrm{NLL}})\leq0.02,
\qquad
\Delta_{\mathcal E}+2\operatorname{MCSE}(\Delta_{\mathcal E})<0.
\label{eq:selection-guard}
\end{equation}
MCSE为蒙特卡洛标准误（Monte Carlo standard error），量化有限数值重复带来的波动。
同一项复核中，两组参数使用配对随机种子；似然复核使用8次独立的1024粒子滤波，
生成复核为每组参数另生成256条训练路径。候选排序、两项复核及最终评价分别使用独立种子组。
标准误通过删除一次完整滤波重复或一条完整路径的刀切法计算，保留差值的配对关系。

每次选参执行一批生成复核。任一条件未通过时保留$\theta_{\mathrm N}$；
两候选相同时直接保留该参数。两项复核完成后冻结参数，再评价后段行为与口述。
这些阈值和预算规定本模型的当前估计协议，独立于个体认知参数。
两倍模拟标准误用于控制数值选择的不确定性，其含义限于模拟误差。

\subsection{条件预测与自主生成各检验什么}
条件预测每题使用此前真实选择与反馈，评价模型在已有经验下对下一题的预测。
后段顺序预测可继续吸收已经发生的后段观测；后段评分与该个体训练参数选择分开。
从初始状态开始的自主生成自行作答、获得相应环境反馈并继续学习，
用以检查规则发现时间、学习形状、停滞、恢复与末段表现的分布。

自主生成保留完整学习路径，并评价路径分布的位置、离散程度、对观察值的覆盖和分布评分。
分布越窄意味着预测越集中；其价值同时取决于中心位置与覆盖是否合理。
个体参数由训练似然约束下的生成选参确定，机制判断结合后段行为、自主分布、内部状态检验和复杂度。
当某条观察学习曲线落在模拟分布边缘时，应继续定位模型过程与估计误差的来源。

\section{口述：对当前作答依据的带噪声检验}
\label{sec:oral}

\subsection{检验对象与时点}
口述指导语要求说明：当前刺激的哪些特点使被试选择这个选项。
因此，主要检验对象为\textbf{本题选择之后、反馈之前的局部作答依据}。
模型仅由行为估参和推断状态；实际口述经编码成为一组判断约束，
再检验这些内容与模型推断的作答依据相容到什么程度。

口述是内部表征的近似检验，报告还受到表达、遗漏、参照歧义和事后解释的影响。
口述检验考察报告内容与推断的学习过程是否相容，并保留不同报告层次的噪声与覆盖限制。
口述质量分层使用独立于模型成绩的标准，常规多部位表达的低目录覆盖单独记录。

\Needspace{10\baselineskip}
\subsection{局部规则分量的解释质量}
设规则$h$对所选类别的区域有局部分量$j$，$\alpha_{t,h,j}^{(n)}$为该粒子内的分量权重。
选择到内部刺激最近的分量，并列时均分；分量权重之和为1。
实际选择之后，粒子、规则和分量的联合解释质量，即分配给各潜在作答依据的概率质量，为
\begin{equation}
A_t(n,h,j\mid y_t)=
\frac{W_t^{(n),\mathrm{pre}}\pi_t^{(n)}(h)p_{t,h}^{(n)}(y_t)
\alpha_{t,h,j}^{(n)}}{P_t(y_t)}.
\label{eq:oral-responsibility}
\end{equation}
该质量对全部$n,h,j$求和为1。
它按当前选择对潜在作答依据条件化一次，把更能解释该回答的规则分量赋予更大权重。

对一条报告的完整约束，结构相容要求全部已编码判断在局部规则中获得显式支持；
功能相容要求整个局部分量区域蕴含这些数值约束。
区域与报告约束没有正体积交集时记为冲突，其他信息不足情形记为未确定。
同一报告的多个判断在同一局部分量内联合评价。
把式\eqref{eq:oral-responsibility}中相应质量累加，分别得到
\begin{equation}
M_t^{\mathrm{compatible}}+M_t^{\mathrm{conflict}}+M_t^{\mathrm{undetermined}}=1.
\end{equation}
这些指标量化内容相容性，报告措辞的生成概率需要另外定义表达模型。

\subsection{编码、覆盖与解释边界}
编码保存原话、解析覆盖、参照歧义和目录范围。
带“一半”的长短按长度范围中点解释；无明确参照的“长／短”具有歧义，
保留预定的中点与相对读法作敏感性分析。
在同一报告集合上使用预定读法与判定标准，并报告相容、冲突、未确定质量及覆盖率。

覆盖率检查目录中是否至少存在一个分量能够支持整条报告，可在模型评价之前确定。
目录覆盖低意味着该对齐方法对相应内容的判别力有限。
多部位报告可能受到目录覆盖的限制，因此报告内容与方法覆盖率分别记录。
对于目录可覆盖内容，进一步比较相容规则是否在场、承担多少作答贡献，才能定位状态或机制偏差。

口述未进入个体参数拟合，部分口述已参与机制开发。
开发样本上的吻合支持模型解释的可行性，独立验证需使用事先冻结的方法及独立资料。
状态变化对口述解释的贡献，可通过同一模型的动态表征与静态表征进行比较。

\section{过程预测与模型检验}
\label{sec:validation}

\subsection{规则发现、保留与使用的过程预测}

规则信念将行为变化与规则管理联系起来。在活跃集合固定时，某条规则的相对信念提高，
会增加其作答权重和作为局部父规则的概率，并提高它在容量替换时的保留优先级。
失败累积量控制搜索概率；规则组成、信念比例和判断精度的变化共同形成学习曲线。

反馈调节机制还给出具体的过程预测。正确反馈保持本次更新前的信念比例，再按反馈似然
区分规则；错误反馈在$\delta<1$时降低已有信念的集中程度，并在发生替换时给予新规则均匀的试用份额。
部分正确反馈产生居中的保持率和入场份额。可分别记录这些状态变化，检验发现新规则、
保持已有规则与实际使用规则之间的联系。

关键过程包括目标规则何时入场、持续活跃多久、何时获得主要作答贡献，
以及它的判断精度与配对信念如何变化。把这些过程与选择概率同时分析，
可以定位行为变化对应的内部状态变化。

\subsection{参数恢复与状态恢复}
参数恢复检验已知生成参数能否从模拟行为中重新估计。
在给定参数下生成完整选择与反馈序列，再按相同估计方法拟合，
比较生成参数与估计参数、两者的行为似然，以及重复模拟中的估计变异。
这一检验同时反映数据量、参数之间的补偿关系和有限搜索对估计的影响。

状态恢复检验已知内部过程能否从行为中重建。
模拟时保留真实规则集合、规则信念、判断精度和配对信念，再仅根据生成的行为序列推断这些状态。
将推断分布与真实路径比较，可以检查规则在场概率、作答贡献和精度轨迹的恢复程度。
参数恢复与状态恢复共同约束个体参数及内部表征的心理解释。

\subsection{模型的适用范围与解释边界}
规则目录规定模型能够表达的判断形式，固定类别标签规定各类别区域对应的反应。
因此，规则发现、容量和口述相容性的解释均以该目录为基础。
口述内容超出目录表达范围时，覆盖率用于说明相关评价的判别能力。

反馈调节保持率与入场份额，规定成功、失败及部分成功如何改变已有判断和新规则的试用程度。
新规则继承加权平均精度，表示任务中形成的刺激区分能力能够在规则之间延续。
条件3中新规则继承综合配对知识，表示已有类别关系知识能够用于后续规则。
这些心理假设的适用性需要结合行为预测、自主生成与内部状态检验评价。

个体参数的心理解释取决于估计可靠性。相近行为似然可能对应不同参数或内部轨迹，
有限粒子数也会带来数值误差。稳定的个体参数解释和独立验证仍需进一步检验。
使用过的开发资料与独立评价资料分别记录；研究结论同时注明参数估计范围、
数值精度、目录覆盖率及所用资料的验证范围。

\appendix
\section{固定设置与初始化}
\label{app:settings}

下表汇总模型的固定设置。个体参数的定义与范围见第\ref{sec:parameters}节。

{\small\noindent
\begin{tabularx}{\linewidth}{L{3.3cm}Y}
\toprule
项目 & 规定\\
\midrule
基础规则权重 & 目录中每条规则的基础权重为$\mu(h)=1/J$。\\
首题状态 & 从基础权重抽取一条规则，其信念为1，判断精度为$\beta_0$，失败累积量为0。\\
首题配对知识 & 条件3的三种配对等概率，各为$1/3$。\\
知觉误差 & 分布形式与参数由独立知觉测量确定，在行为参数估计期间固定。\\
判断精度 & 取值限制为$[0.1,25]$，反馈梯度的更新步长为$10\eta$。\\
额外随机作答率 & 固定为0；选择概率由活跃规则的加权预测给出。\\
局部结构搜索 & 在非空结构邻域中，按编辑距离的负指数归一化抽取候选。\\
局部功能搜索 & 结构邻域为空时使用功能距离，指数尺度为$0.10$；距离以100000个均匀参考刺激估计，随机种子为0。\\
规则信念保持 & 有效保持率为$\delta_t=\delta+(1-\delta)r_t$；幂变换并归一化后结合反馈似然。\\
容量替换 & 新规则替换更新后信念最低的已有规则；并列时等概率选择。\\
新规则信念 & 有空位时获得$1/(|\A_t|+1)$；满容量时获得$r_t\pi_t^+(h^-)+(1-r_t)/K$。幸存规则按原比例分配剩余质量。\\
新规则继承 & 判断精度继承替换前全部活跃规则的信念加权平均值；条件3继承综合配对知识。\\
\bottomrule
\end{tabularx}}

\section{规则目录}
\label{app:catalogue}

\subsection{二分类目录：29条}
{\small\noindent
\begin{tabularx}{\linewidth}{L{4.4cm}Yr}
\toprule
规则形式 & 判断例子或边界 & 数量\\
\midrule
单特征阈值 & $x_i\leq0.5$ & 4\\
两个特征的顺序 & $x_i\leq x_j$ & 6\\
两个特征的和 & $x_i+x_j\leq1$ & 6\\
两组特征和的顺序 & $x_i+x_j\leq x_k+x_l$ & 3\\
两个特征是否接近 & $|x_i-x_j|\leq0.10$ & 6\\
单特征是否在中间 & $|x_i-0.5|\leq0.10$ & 4\\
\midrule
合计 & 类别区域与标签固定 & 29\\
\bottomrule
\end{tabularx}}

\Needspace{18\baselineskip}
\subsection{四分类目录：152条}
{\small
\begin{longtable}{L{2.6cm}L{11.1cm}r}
\toprule
涉及特征数 & 规则形式 & 数量\\
\midrule\endhead
两个特征 & 两个单特征阈值交叉 & 6\\
两个特征 & 同一特征对的顺序与和交叉 & 6\\
三个特征 & 阈值与另两特征顺序 & 12\\
三个特征 & 阈值与另两特征和 & 12\\
三个特征 & 三阈值条件分支 & 24\\
三个特征 & 阈值分支下顺序或和 & 24\\
四个特征 & 不重叠特征对的顺序交叉 & 3\\
四个特征 & 不重叠特征对的和交叉 & 3\\
四个特征 & 顺序分支与剩余阈值 & 12\\
四个特征 & 和分支与剩余阈值 & 12\\
四个特征 & 最大或最小特征分类 & 2\\
两个特征 & 阈值与共享特征顺序 & 12\\
三个特征 & 围绕同一部位两次比较 & 12\\
三个特征 & 两部位接近与第三部位长短交叉 & 12\\
\midrule
合计 & 固定标签下的规则判断与类别区域 & 152\\
\bottomrule
\end{longtable}}

\clearpage
\section{主要符号}
\label{app:symbols}

{\small
\begin{longtable}{L{3.0cm}L{11.5cm}}
\toprule
符号 & 含义\\
\midrule\endhead
$S_t,\ m_t$ & 决策就绪状态与本题内部记录。\\
$\A_t,\ \Hh,\ K$ & 活跃规则集合、固定规则目录与容量。\\

$\pi_t(h),\ \bar\pi_t(h),\ \pi_t^+(h)$ & 作答前规则信念、保持处理后的中间信念、反馈更新后的规则信念。\\
$\delta,\ \delta_t$ & 基础规则信念保持率与由本题反馈决定的有效保持率。\\
$\nu_{t+1},\ \A_t^{\mathrm{keep}}$ & 新规则的入场信念份额与本次规则管理后的幸存集合。\\
$\beta_t(h),\ \beta_t^+(h),\ \bar\beta_t$ & 作答时精度、反馈学习后精度与新规则继承的平均精度。\\
$F_t$ & 决策就绪状态中的失败累积量。\\
$q_t,\ \bar q_t,\ q_t^+,\ \gamma$ & 条件配对信念、保持后分布、反馈后分布与配对保持率。\\
$\Pi_t,\ \Pi_t^+,\ Q_t^+$ & 作答前及反馈后的规则与配对联合信念、反馈后的整体配对边际。\\
$\ell_t(h),\ \ell_t(h,k)$ & 规则的反馈似然与给定规则和配对时的反馈似然。\\
$p_{t,h}(c),\ p_t(c)$ & 单条规则的类别预测与单条路径的混合选择预测。\\
$P_t(c),\ W_t^{(n)}$ & 研究者跨路径的预测概率与粒子权重。\\
$\Dd,\ \Ll,\ k_\theta$ & 决策核、学习核与单试次联合生成核。\\
$\Ii,\ \Tt$ & 轨迹推断算子与状态演化算子。\\
$\phi,\ \widehat{\mathcal E}$ & 自主路径的行为特征向量及其Energy评分。\\
$\theta_{\mathrm N},\ \theta_{\mathrm G}$ & 训练似然候选与训练生成提议参数。\\

\bottomrule
\end{longtable}}
\end{document}
